%%%%%%%%%%%%%%%%%%%%%%%%%%%%%%%%%%%%%%%%%%%%%%%%%%%%%%%%%%%%
% 14.9 DERIVATIVE TABLES
% The Composition of Advanced Mathematics
%%%%%%%%%%%%%%%%%%%%%%%%%%%%%%%%%%%%%%%%%%%%%%%%%%%%%%%%%%%%

\section{Derivative Tables}
\label{sec:derivative-tables}

\prerequisite{
Functions, limits, algebraic manipulation, trigonometric functions,
inverse trigonometric functions, exponential and logarithmic functions,
hyperbolic functions, implicit differentiation, and multivariable notation.
}

\learningobjectives{
By the end of this reference section, the reader should be able to:
\begin{itemize}
    \item recognize standard derivative notation;
    \item use constant, power, sum, difference, product, quotient,
          and chain rules;
    \item differentiate polynomial and rational functions;
    \item differentiate exponential and logarithmic functions;
    \item differentiate all six trigonometric functions;
    \item differentiate inverse trigonometric functions;
    \item differentiate hyperbolic and inverse hyperbolic functions;
    \item differentiate functions involving absolute value where appropriate;
    \item use logarithmic differentiation;
    \item perform implicit differentiation;
    \item compute higher derivatives;
    \item recognize derivatives of common composite functions;
    \item differentiate parametric equations;
    \item compute derivatives in polar form;
    \item recognize partial derivatives and gradients;
    \item compute Jacobians and Hessians;
    \item recognize directional derivatives;
    \item differentiate vector-valued functions;
    \item identify common derivative patterns used in differential equations,
          optimization, probability, and applied mathematics;
    \item avoid common sign, chain-rule, and domain errors.
\end{itemize}
}

%%%%%%%%%%%%%%%%%%%%%%%%%%%%%%%%%%%%%%%%%%%%%%%%%%%%%%%%%%%%
\subsection{Derivative Notation}
\label{subsec:derivative-notation-table-reference}
%%%%%%%%%%%%%%%%%%%%%%%%%%%%%%%%%%%%%%%%%%%%%%%%%%%%%%%%%%%%

For

\[
y=f(x),
\]

common derivative notation includes

\[
\boxed{
f'(x),
}
\]

\[
\boxed{
\frac{dy}{dx},
}
\]

\[
\boxed{
\frac{df}{dx},
}
\]

and

\[
\boxed{
D_xf.
}
\]

All represent the derivative with respect to \(x\).

%%%%%%%%%%%%%%%%%%%%%%%%%%%%%%%%%%%%%%%%%%%%%%%%%%%%%%%%%%%%
\subsection{Derivative Definition}
\label{subsec:derivative-definition-reference}
%%%%%%%%%%%%%%%%%%%%%%%%%%%%%%%%%%%%%%%%%%%%%%%%%%%%%%%%%%%%

The derivative of \(f\) at \(x\) is

\[
\boxed{
f'(x)
=
\lim_{h\to0}
\frac{
f(x+h)-f(x)
}{
h
}
}
\]

whenever the limit exists.

An equivalent form at \(x=a\) is

\[
\boxed{
f'(a)
=
\lim_{x\to a}
\frac{
f(x)-f(a)
}{
x-a
}.
}
\]

%%%%%%%%%%%%%%%%%%%%%%%%%%%%%%%%%%%%%%%%%%%%%%%%%%%%%%%%%%%%
\subsection{Derivative as Rate of Change}
\label{subsec:derivative-rate-change-reference}
%%%%%%%%%%%%%%%%%%%%%%%%%%%%%%%%%%%%%%%%%%%%%%%%%%%%%%%%%%%%

The derivative represents instantaneous rate of change.

If

\[
y=f(x),
\]

then

\[
\boxed{
f'(a)
}
\]

is the slope of the tangent line at

\[
x=a.
\]

%%%%%%%%%%%%%%%%%%%%%%%%%%%%%%%%%%%%%%%%%%%%%%%%%%%%%%%%%%%%
\subsection{Constant Rule}
\label{subsec:constant-derivative-rule}
%%%%%%%%%%%%%%%%%%%%%%%%%%%%%%%%%%%%%%%%%%%%%%%%%%%%%%%%%%%%

For constant \(c\),

\[
\boxed{
\frac{d}{dx}c
=
0.
}
\]

%%%%%%%%%%%%%%%%%%%%%%%%%%%%%%%%%%%%%%%%%%%%%%%%%%%%%%%%%%%%
\subsection{Power Rule}
\label{subsec:power-derivative-rule}
%%%%%%%%%%%%%%%%%%%%%%%%%%%%%%%%%%%%%%%%%%%%%%%%%%%%%%%%%%%%

For real exponent \(n\), wherever the expression is differentiable,

\[
\boxed{
\frac{d}{dx}x^n
=
nx^{n-1}.
}
\]

For integer \(n\), this is the familiar power rule.

%%%%%%%%%%%%%%%%%%%%%%%%%%%%%%%%%%%%%%%%%%%%%%%%%%%%%%%%%%%%
\subsection{Constant Multiple Rule}
\label{subsec:constant-multiple-derivative-rule}
%%%%%%%%%%%%%%%%%%%%%%%%%%%%%%%%%%%%%%%%%%%%%%%%%%%%%%%%%%%%

For constant \(c\),

\[
\boxed{
\frac{d}{dx}
[c f(x)]
=
c f'(x).
}
\]

%%%%%%%%%%%%%%%%%%%%%%%%%%%%%%%%%%%%%%%%%%%%%%%%%%%%%%%%%%%%
\subsection{Sum Rule}
\label{subsec:sum-derivative-rule}
%%%%%%%%%%%%%%%%%%%%%%%%%%%%%%%%%%%%%%%%%%%%%%%%%%%%%%%%%%%%

\[
\boxed{
\frac{d}{dx}
[f(x)+g(x)]
=
f'(x)+g'(x).
}
\]

%%%%%%%%%%%%%%%%%%%%%%%%%%%%%%%%%%%%%%%%%%%%%%%%%%%%%%%%%%%%
\subsection{Difference Rule}
\label{subsec:difference-derivative-rule}
%%%%%%%%%%%%%%%%%%%%%%%%%%%%%%%%%%%%%%%%%%%%%%%%%%%%%%%%%%%%

\[
\boxed{
\frac{d}{dx}
[f(x)-g(x)]
=
f'(x)-g'(x).
}
\]

%%%%%%%%%%%%%%%%%%%%%%%%%%%%%%%%%%%%%%%%%%%%%%%%%%%%%%%%%%%%
\subsection{Product Rule}
\label{subsec:product-derivative-rule}
%%%%%%%%%%%%%%%%%%%%%%%%%%%%%%%%%%%%%%%%%%%%%%%%%%%%%%%%%%%%

\[
\boxed{
\frac{d}{dx}
[f(x)g(x)]
=
f'(x)g(x)
+
f(x)g'(x).
}
\]

A useful mnemonic is

\[
\boxed{
(fg)'
=
f'g+fg'.
}
\]

%%%%%%%%%%%%%%%%%%%%%%%%%%%%%%%%%%%%%%%%%%%%%%%%%%%%%%%%%%%%
\subsection{Quotient Rule}
\label{subsec:quotient-derivative-rule}
%%%%%%%%%%%%%%%%%%%%%%%%%%%%%%%%%%%%%%%%%%%%%%%%%%%%%%%%%%%%

For

\[
g(x)\neq0,
\]

\[
\boxed{
\frac{d}{dx}
\left(
\frac{f(x)}{g(x)}
\right)
=
\frac{
f'(x)g(x)-f(x)g'(x)
}{
[g(x)]^2
}.
}
\]

%%%%%%%%%%%%%%%%%%%%%%%%%%%%%%%%%%%%%%%%%%%%%%%%%%%%%%%%%%%%
\subsection{Chain Rule}
\label{subsec:chain-derivative-rule}
%%%%%%%%%%%%%%%%%%%%%%%%%%%%%%%%%%%%%%%%%%%%%%%%%%%%%%%%%%%%

For

\[
y=f(g(x)),
\]

\[
\boxed{
\frac{dy}{dx}
=
f'(g(x))g'(x).
}
\]

In Leibniz notation,

\[
\boxed{
\frac{dy}{dx}
=
\frac{dy}{du}
\frac{du}{dx}.
}
\]

%%%%%%%%%%%%%%%%%%%%%%%%%%%%%%%%%%%%%%%%%%%%%%%%%%%%%%%%%%%%
\subsection{Worked Example: Chain Rule}
\label{subsec:worked-chain-rule-reference}
%%%%%%%%%%%%%%%%%%%%%%%%%%%%%%%%%%%%%%%%%%%%%%%%%%%%%%%%%%%%

\begin{workedexamplebox}{Differentiate a Composite Power}

Differentiate

\[
y
=
\left(
(x+1)^2+3
\right)^4.
\]

Let

\[
u
=
(x+1)^2+3.
\]

Then

\[
y=u^4.
\]

Differentiate:

\[
\frac{dy}{du}
=
4u^3,
\]

and

\[
\frac{du}{dx}
=
2(x+1).
\]

Therefore,

\begin{answerbox}
\[
\boxed{
y'
=
8(x+1)
\left(
(x+1)^2+3
\right)^3.
}
\]
\end{answerbox}

\end{workedexamplebox}

%%%%%%%%%%%%%%%%%%%%%%%%%%%%%%%%%%%%%%%%%%%%%%%%%%%%%%%%%%%%
\subsection{General Power-Chain Rule}
\label{subsec:power-chain-rule}
%%%%%%%%%%%%%%%%%%%%%%%%%%%%%%%%%%%%%%%%%%%%%%%%%%%%%%%%%%%%

For differentiable \(u=u(x)\),

\[
\boxed{
\frac{d}{dx}
[u(x)]^n
=
n[u(x)]^{n-1}u'(x).
}
\]

%%%%%%%%%%%%%%%%%%%%%%%%%%%%%%%%%%%%%%%%%%%%%%%%%%%%%%%%%%%%
\subsection{Polynomial Derivatives}
\label{subsec:polynomial-derivatives-reference}
%%%%%%%%%%%%%%%%%%%%%%%%%%%%%%%%%%%%%%%%%%%%%%%%%%%%%%%%%%%%

For

\[
p(x)
=
a_nx^n+\cdots+a_1x+a_0,
\]

\[
\boxed{
p'(x)
=
na_nx^{n-1}
+
\cdots
+
a_1.
}
\]

%%%%%%%%%%%%%%%%%%%%%%%%%%%%%%%%%%%%%%%%%%%%%%%%%%%%%%%%%%%%
\subsection{Reciprocal Function}
\label{subsec:reciprocal-derivative-reference}
%%%%%%%%%%%%%%%%%%%%%%%%%%%%%%%%%%%%%%%%%%%%%%%%%%%%%%%%%%%%

For

\[
x\neq0,
\]

\[
\boxed{
\frac{d}{dx}
\left(
\frac1x
\right)
=
-\frac1{x^2}.
}
\]

More generally,

\[
\boxed{
\frac{d}{dx}
\left(
\frac1{u}
\right)
=
-\frac{u'}{u^2}.
}
\]

%%%%%%%%%%%%%%%%%%%%%%%%%%%%%%%%%%%%%%%%%%%%%%%%%%%%%%%%%%%%
\subsection{Square Root}
\label{subsec:square-root-derivative-reference}
%%%%%%%%%%%%%%%%%%%%%%%%%%%%%%%%%%%%%%%%%%%%%%%%%%%%%%%%%%%%

For

\[
x>0,
\]

\[
\boxed{
\frac{d}{dx}
\sqrt{x}
=
\frac{1}{2\sqrt{x}}.
}
\]

For a differentiable inner function \(u(x)>0\),

\[
\boxed{
\frac{d}{dx}
\sqrt{u}
=
\frac{u'}{2\sqrt{u}}.
}
\]

%%%%%%%%%%%%%%%%%%%%%%%%%%%%%%%%%%%%%%%%%%%%%%%%%%%%%%%%%%%%
\subsection{General Radical}
\label{subsec:general-radical-derivative}
%%%%%%%%%%%%%%%%%%%%%%%%%%%%%%%%%%%%%%%%%%%%%%%%%%%%%%%%%%%%

For suitable domains,

\[
\boxed{
\frac{d}{dx}
u^{1/n}
=
\frac1n
u^{1/n-1}
u'.
}
\]

%%%%%%%%%%%%%%%%%%%%%%%%%%%%%%%%%%%%%%%%%%%%%%%%%%%%%%%%%%%%
\subsection{Natural Exponential}
\label{subsec:e-exponential-derivative}
%%%%%%%%%%%%%%%%%%%%%%%%%%%%%%%%%%%%%%%%%%%%%%%%%%%%%%%%%%%%

\[
\boxed{
\frac{d}{dx}e^x
=
e^x.
}
\]

For \(u=u(x)\),

\[
\boxed{
\frac{d}{dx}e^u
=
e^u u'.
}
\]

%%%%%%%%%%%%%%%%%%%%%%%%%%%%%%%%%%%%%%%%%%%%%%%%%%%%%%%%%%%%
\subsection{General Exponential}
\label{subsec:general-exponential-derivative}
%%%%%%%%%%%%%%%%%%%%%%%%%%%%%%%%%%%%%%%%%%%%%%%%%%%%%%%%%%%%

For constant

\[
a>0,
\qquad
a\neq1,
\]

\[
\boxed{
\frac{d}{dx}a^x
=
a^x\ln a.
}
\]

For \(u=u(x)\),

\[
\boxed{
\frac{d}{dx}
a^{u}
=
a^u\ln(a)\,u'.
}
\]

%%%%%%%%%%%%%%%%%%%%%%%%%%%%%%%%%%%%%%%%%%%%%%%%%%%%%%%%%%%%
\subsection{Natural Logarithm}
\label{subsec:natural-logarithm-derivative}
%%%%%%%%%%%%%%%%%%%%%%%%%%%%%%%%%%%%%%%%%%%%%%%%%%%%%%%%%%%%

For

\[
x>0,
\]

\[
\boxed{
\frac{d}{dx}\ln x
=
\frac1x.
}
\]

More generally, where \(u(x)\neq0\),

\[
\boxed{
\frac{d}{dx}
\ln|u|
=
\frac{u'}{u}.
}
\]

%%%%%%%%%%%%%%%%%%%%%%%%%%%%%%%%%%%%%%%%%%%%%%%%%%%%%%%%%%%%
\subsection{General Logarithm}
\label{subsec:general-logarithm-derivative}
%%%%%%%%%%%%%%%%%%%%%%%%%%%%%%%%%%%%%%%%%%%%%%%%%%%%%%%%%%%%

For

\[
a>0,
\qquad
a\neq1,
\qquad
x>0,
\]

\[
\boxed{
\frac{d}{dx}
\log_a x
=
\frac{
1
}{
x\ln a
}.
}
\]

For \(u=u(x)>0\),

\[
\boxed{
\frac{d}{dx}
\log_a u
=
\frac{
u'
}{
u\ln a
}.
}
\]

%%%%%%%%%%%%%%%%%%%%%%%%%%%%%%%%%%%%%%%%%%%%%%%%%%%%%%%%%%%%
\subsection{Sine Derivative}
\label{subsec:sine-derivative-reference}
%%%%%%%%%%%%%%%%%%%%%%%%%%%%%%%%%%%%%%%%%%%%%%%%%%%%%%%%%%%%

\[
\boxed{
\frac{d}{dx}\sin x
=
\cos x.
}
\]

For \(u=u(x)\),

\[
\boxed{
\frac{d}{dx}\sin u
=
\cos u\,u'.
}
\]

%%%%%%%%%%%%%%%%%%%%%%%%%%%%%%%%%%%%%%%%%%%%%%%%%%%%%%%%%%%%
\subsection{Cosine Derivative}
\label{subsec:cosine-derivative-reference}
%%%%%%%%%%%%%%%%%%%%%%%%%%%%%%%%%%%%%%%%%%%%%%%%%%%%%%%%%%%%

\[
\boxed{
\frac{d}{dx}\cos x
=
-\sin x.
}
\]

For \(u=u(x)\),

\[
\boxed{
\frac{d}{dx}\cos u
=
-\sin u\,u'.
}
\]

%%%%%%%%%%%%%%%%%%%%%%%%%%%%%%%%%%%%%%%%%%%%%%%%%%%%%%%%%%%%
\subsection{Tangent Derivative}
\label{subsec:tangent-derivative-reference}
%%%%%%%%%%%%%%%%%%%%%%%%%%%%%%%%%%%%%%%%%%%%%%%%%%%%%%%%%%%%

\[
\boxed{
\frac{d}{dx}\tan x
=
\sec^2x.
}
\]

For \(u=u(x)\),

\[
\boxed{
\frac{d}{dx}\tan u
=
\sec^2u\,u'.
}
\]

%%%%%%%%%%%%%%%%%%%%%%%%%%%%%%%%%%%%%%%%%%%%%%%%%%%%%%%%%%%%
\subsection{Cotangent Derivative}
\label{subsec:cotangent-derivative-reference}
%%%%%%%%%%%%%%%%%%%%%%%%%%%%%%%%%%%%%%%%%%%%%%%%%%%%%%%%%%%%

\[
\boxed{
\frac{d}{dx}\cot x
=
-\csc^2x.
}
\]

For \(u=u(x)\),

\[
\boxed{
\frac{d}{dx}\cot u
=
-\csc^2u\,u'.
}
\]

%%%%%%%%%%%%%%%%%%%%%%%%%%%%%%%%%%%%%%%%%%%%%%%%%%%%%%%%%%%%
\subsection{Secant Derivative}
\label{subsec:secant-derivative-reference}
%%%%%%%%%%%%%%%%%%%%%%%%%%%%%%%%%%%%%%%%%%%%%%%%%%%%%%%%%%%%

\[
\boxed{
\frac{d}{dx}\sec x
=
\sec x\tan x.
}
\]

For \(u=u(x)\),

\[
\boxed{
\frac{d}{dx}\sec u
=
\sec u\tan u\,u'.
}
\]

%%%%%%%%%%%%%%%%%%%%%%%%%%%%%%%%%%%%%%%%%%%%%%%%%%%%%%%%%%%%
\subsection{Cosecant Derivative}
\label{subsec:cosecant-derivative-reference}
%%%%%%%%%%%%%%%%%%%%%%%%%%%%%%%%%%%%%%%%%%%%%%%%%%%%%%%%%%%%

\[
\boxed{
\frac{d}{dx}\csc x
=
-\csc x\cot x.
}
\]

For \(u=u(x)\),

\[
\boxed{
\frac{d}{dx}\csc u
=
-\csc u\cot u\,u'.
}
\]

%%%%%%%%%%%%%%%%%%%%%%%%%%%%%%%%%%%%%%%%%%%%%%%%%%%%%%%%%%%%
\subsection{Quick Trigonometric Derivative Table}
\label{subsec:quick-trig-derivative-table}
%%%%%%%%%%%%%%%%%%%%%%%%%%%%%%%%%%%%%%%%%%%%%%%%%%%%%%%%%%%%

\[
\boxed{
(\sin x)'=\cos x
}
\]

\[
\boxed{
(\cos x)'=-\sin x
}
\]

\[
\boxed{
(\tan x)'=\sec^2x
}
\]

\[
\boxed{
(\cot x)'=-\csc^2x
}
\]

\[
\boxed{
(\sec x)'=\sec x\tan x
}
\]

\[
\boxed{
(\csc x)'=-\csc x\cot x
}
\]

%%%%%%%%%%%%%%%%%%%%%%%%%%%%%%%%%%%%%%%%%%%%%%%%%%%%%%%%%%%%
\subsection{Inverse Sine Derivative}
\label{subsec:inverse-sine-derivative}
%%%%%%%%%%%%%%%%%%%%%%%%%%%%%%%%%%%%%%%%%%%%%%%%%%%%%%%%%%%%

For

\[
|x|<1,
\]

\[
\boxed{
\frac{d}{dx}\arcsin x
=
\frac{1}{\sqrt{1-x^2}}.
}
\]

For \(u=u(x)\),

\[
\boxed{
\frac{d}{dx}\arcsin u
=
\frac{
u'
}{
\sqrt{1-u^2}
}.
}
\]

%%%%%%%%%%%%%%%%%%%%%%%%%%%%%%%%%%%%%%%%%%%%%%%%%%%%%%%%%%%%
\subsection{Inverse Cosine Derivative}
\label{subsec:inverse-cosine-derivative}
%%%%%%%%%%%%%%%%%%%%%%%%%%%%%%%%%%%%%%%%%%%%%%%%%%%%%%%%%%%%

For

\[
|x|<1,
\]

\[
\boxed{
\frac{d}{dx}\arccos x
=
-\frac{1}{\sqrt{1-x^2}}.
}
\]

For \(u=u(x)\),

\[
\boxed{
\frac{d}{dx}\arccos u
=
-\frac{
u'
}{
\sqrt{1-u^2}
}.
}
\]

%%%%%%%%%%%%%%%%%%%%%%%%%%%%%%%%%%%%%%%%%%%%%%%%%%%%%%%%%%%%
\subsection{Inverse Tangent Derivative}
\label{subsec:inverse-tangent-derivative}
%%%%%%%%%%%%%%%%%%%%%%%%%%%%%%%%%%%%%%%%%%%%%%%%%%%%%%%%%%%%

\[
\boxed{
\frac{d}{dx}\arctan x
=
\frac{1}{1+x^2}.
}
\]

For \(u=u(x)\),

\[
\boxed{
\frac{d}{dx}\arctan u
=
\frac{
u'
}{
1+u^2
}.
}
\]

%%%%%%%%%%%%%%%%%%%%%%%%%%%%%%%%%%%%%%%%%%%%%%%%%%%%%%%%%%%%
\subsection{Inverse Cotangent Derivative}
\label{subsec:inverse-cotangent-derivative}
%%%%%%%%%%%%%%%%%%%%%%%%%%%%%%%%%%%%%%%%%%%%%%%%%%%%%%%%%%%%

Under the common principal-value convention,

\[
\boxed{
\frac{d}{dx}\operatorname{arccot}x
=
-\frac{1}{1+x^2}.
}
\]

For \(u=u(x)\),

\[
\boxed{
\frac{d}{dx}\operatorname{arccot}u
=
-\frac{
u'
}{
1+u^2
}.
}
\]

Inverse-cotangent conventions vary between texts.

%%%%%%%%%%%%%%%%%%%%%%%%%%%%%%%%%%%%%%%%%%%%%%%%%%%%%%%%%%%%
\subsection{Inverse Secant Derivative}
\label{subsec:inverse-secant-derivative}
%%%%%%%%%%%%%%%%%%%%%%%%%%%%%%%%%%%%%%%%%%%%%%%%%%%%%%%%%%%%

For

\[
|x|>1,
\]

a standard real derivative formula is

\[
\boxed{
\frac{d}{dx}\operatorname{arcsec}x
=
\frac{
1
}{
|x|\sqrt{x^2-1}
}.
}
\]

For \(u=u(x)\),

\[
\boxed{
\frac{d}{dx}\operatorname{arcsec}u
=
\frac{
u'
}{
|u|\sqrt{u^2-1}
}.
}
\]

%%%%%%%%%%%%%%%%%%%%%%%%%%%%%%%%%%%%%%%%%%%%%%%%%%%%%%%%%%%%
\subsection{Inverse Cosecant Derivative}
\label{subsec:inverse-cosecant-derivative}
%%%%%%%%%%%%%%%%%%%%%%%%%%%%%%%%%%%%%%%%%%%%%%%%%%%%%%%%%%%%

For

\[
|x|>1,
\]

\[
\boxed{
\frac{d}{dx}\operatorname{arccsc}x
=
-\frac{
1
}{
|x|\sqrt{x^2-1}
}.
}
\]

For \(u=u(x)\),

\[
\boxed{
\frac{d}{dx}\operatorname{arccsc}u
=
-\frac{
u'
}{
|u|\sqrt{u^2-1}
}.
}
\]

%%%%%%%%%%%%%%%%%%%%%%%%%%%%%%%%%%%%%%%%%%%%%%%%%%%%%%%%%%%%
\subsection{Quick Inverse-Trigonometric Derivative Table}
\label{subsec:quick-inverse-trig-derivative-table}
%%%%%%%%%%%%%%%%%%%%%%%%%%%%%%%%%%%%%%%%%%%%%%%%%%%%%%%%%%%%

\[
\boxed{
(\arcsin x)'
=
\frac{1}{\sqrt{1-x^2}}
}
\]

\[
\boxed{
(\arccos x)'
=
-\frac{1}{\sqrt{1-x^2}}
}
\]

\[
\boxed{
(\arctan x)'
=
\frac{1}{1+x^2}
}
\]

\[
\boxed{
(\operatorname{arccot}x)'
=
-\frac{1}{1+x^2}
}
\]

\[
\boxed{
(\operatorname{arcsec}x)'
=
\frac{
1
}{
|x|\sqrt{x^2-1}
}
}
\]

\[
\boxed{
(\operatorname{arccsc}x)'
=
-\frac{
1
}{
|x|\sqrt{x^2-1}
}
}
\]

%%%%%%%%%%%%%%%%%%%%%%%%%%%%%%%%%%%%%%%%%%%%%%%%%%%%%%%%%%%%
\subsection{Hyperbolic Sine}
\label{subsec:sinh-derivative-reference}
%%%%%%%%%%%%%%%%%%%%%%%%%%%%%%%%%%%%%%%%%%%%%%%%%%%%%%%%%%%%

\[
\boxed{
\frac{d}{dx}\sinh x
=
\cosh x.
}
\]

For \(u=u(x)\),

\[
\boxed{
\frac{d}{dx}\sinh u
=
\cosh u\,u'.
}
\]

%%%%%%%%%%%%%%%%%%%%%%%%%%%%%%%%%%%%%%%%%%%%%%%%%%%%%%%%%%%%
\subsection{Hyperbolic Cosine}
\label{subsec:cosh-derivative-reference}
%%%%%%%%%%%%%%%%%%%%%%%%%%%%%%%%%%%%%%%%%%%%%%%%%%%%%%%%%%%%

\[
\boxed{
\frac{d}{dx}\cosh x
=
\sinh x.
}
\]

For \(u=u(x)\),

\[
\boxed{
\frac{d}{dx}\cosh u
=
\sinh u\,u'.
}
\]

%%%%%%%%%%%%%%%%%%%%%%%%%%%%%%%%%%%%%%%%%%%%%%%%%%%%%%%%%%%%
\subsection{Hyperbolic Tangent}
\label{subsec:tanh-derivative-reference}
%%%%%%%%%%%%%%%%%%%%%%%%%%%%%%%%%%%%%%%%%%%%%%%%%%%%%%%%%%%%

\[
\boxed{
\frac{d}{dx}\tanh x
=
\operatorname{sech}^2x.
}
\]

For \(u=u(x)\),

\[
\boxed{
\frac{d}{dx}\tanh u
=
\operatorname{sech}^2u\,u'.
}
\]

%%%%%%%%%%%%%%%%%%%%%%%%%%%%%%%%%%%%%%%%%%%%%%%%%%%%%%%%%%%%
\subsection{Hyperbolic Cotangent}
\label{subsec:coth-derivative-reference}
%%%%%%%%%%%%%%%%%%%%%%%%%%%%%%%%%%%%%%%%%%%%%%%%%%%%%%%%%%%%

\[
\boxed{
\frac{d}{dx}\coth x
=
-\operatorname{csch}^2x.
}
\]

%%%%%%%%%%%%%%%%%%%%%%%%%%%%%%%%%%%%%%%%%%%%%%%%%%%%%%%%%%%%
\subsection{Hyperbolic Secant}
\label{subsec:sech-derivative-reference}
%%%%%%%%%%%%%%%%%%%%%%%%%%%%%%%%%%%%%%%%%%%%%%%%%%%%%%%%%%%%

\[
\boxed{
\frac{d}{dx}\operatorname{sech}x
=
-\operatorname{sech}x\tanh x.
}
\]

%%%%%%%%%%%%%%%%%%%%%%%%%%%%%%%%%%%%%%%%%%%%%%%%%%%%%%%%%%%%
\subsection{Hyperbolic Cosecant}
\label{subsec:csch-derivative-reference}
%%%%%%%%%%%%%%%%%%%%%%%%%%%%%%%%%%%%%%%%%%%%%%%%%%%%%%%%%%%%

\[
\boxed{
\frac{d}{dx}\operatorname{csch}x
=
-\operatorname{csch}x\coth x.
}
\]

%%%%%%%%%%%%%%%%%%%%%%%%%%%%%%%%%%%%%%%%%%%%%%%%%%%%%%%%%%%%
\subsection{Quick Hyperbolic Derivative Table}
\label{subsec:quick-hyperbolic-derivative-table}
%%%%%%%%%%%%%%%%%%%%%%%%%%%%%%%%%%%%%%%%%%%%%%%%%%%%%%%%%%%%

\[
\boxed{
(\sinh x)'=\cosh x
}
\]

\[
\boxed{
(\cosh x)'=\sinh x
}
\]

\[
\boxed{
(\tanh x)'=\operatorname{sech}^2x
}
\]

\[
\boxed{
(\coth x)'=-\operatorname{csch}^2x
}
\]

\[
\boxed{
(\operatorname{sech}x)'
=
-\operatorname{sech}x\tanh x
}
\]

\[
\boxed{
(\operatorname{csch}x)'
=
-\operatorname{csch}x\coth x
}
\]

%%%%%%%%%%%%%%%%%%%%%%%%%%%%%%%%%%%%%%%%%%%%%%%%%%%%%%%%%%%%
\subsection{Inverse Hyperbolic Sine}
\label{subsec:inverse-sinh-derivative}
%%%%%%%%%%%%%%%%%%%%%%%%%%%%%%%%%%%%%%%%%%%%%%%%%%%%%%%%%%%%

\[
\boxed{
\frac{d}{dx}
\operatorname{arsinh}x
=
\frac{
1
}{
\sqrt{x^2+1}
}.
}
\]

%%%%%%%%%%%%%%%%%%%%%%%%%%%%%%%%%%%%%%%%%%%%%%%%%%%%%%%%%%%%
\subsection{Inverse Hyperbolic Cosine}
\label{subsec:inverse-cosh-derivative}
%%%%%%%%%%%%%%%%%%%%%%%%%%%%%%%%%%%%%%%%%%%%%%%%%%%%%%%%%%%%

For \(x>1\),

\[
\boxed{
\frac{d}{dx}
\operatorname{arcosh}x
=
\frac{
1
}{
\sqrt{x^2-1}
}.
}
\]

Equivalently, one often writes

\[
\boxed{
\frac{1}{
\sqrt{x-1}\sqrt{x+1}
}.
}
\]

%%%%%%%%%%%%%%%%%%%%%%%%%%%%%%%%%%%%%%%%%%%%%%%%%%%%%%%%%%%%
\subsection{Inverse Hyperbolic Tangent}
\label{subsec:inverse-tanh-derivative}
%%%%%%%%%%%%%%%%%%%%%%%%%%%%%%%%%%%%%%%%%%%%%%%%%%%%%%%%%%%%

For

\[
|x|<1,
\]

\[
\boxed{
\frac{d}{dx}
\operatorname{artanh}x
=
\frac{1}{1-x^2}.
}
\]

%%%%%%%%%%%%%%%%%%%%%%%%%%%%%%%%%%%%%%%%%%%%%%%%%%%%%%%%%%%%
\subsection{Inverse Hyperbolic Cotangent}
\label{subsec:inverse-coth-derivative}
%%%%%%%%%%%%%%%%%%%%%%%%%%%%%%%%%%%%%%%%%%%%%%%%%%%%%%%%%%%%

On standard real domains,

\[
\boxed{
\frac{d}{dx}
\operatorname{arcoth}x
=
\frac{1}{1-x^2}.
}
\]

%%%%%%%%%%%%%%%%%%%%%%%%%%%%%%%%%%%%%%%%%%%%%%%%%%%%%%%%%%%%
\subsection{Inverse Hyperbolic Secant}
\label{subsec:inverse-sech-derivative}
%%%%%%%%%%%%%%%%%%%%%%%%%%%%%%%%%%%%%%%%%%%%%%%%%%%%%%%%%%%%

For

\[
0<x<1,
\]

\[
\boxed{
\frac{d}{dx}
\operatorname{arsech}x
=
-\frac{
1
}{
x\sqrt{1-x^2}
}.
}
\]

%%%%%%%%%%%%%%%%%%%%%%%%%%%%%%%%%%%%%%%%%%%%%%%%%%%%%%%%%%%%
\subsection{Inverse Hyperbolic Cosecant}
\label{subsec:inverse-csch-derivative}
%%%%%%%%%%%%%%%%%%%%%%%%%%%%%%%%%%%%%%%%%%%%%%%%%%%%%%%%%%%%

For \(x\neq0\),

\[
\boxed{
\frac{d}{dx}
\operatorname{arcsch}x
=
-\frac{
1
}{
|x|\sqrt{1+x^2}
}.
}
\]

Notation and branch conventions can vary between texts.

%%%%%%%%%%%%%%%%%%%%%%%%%%%%%%%%%%%%%%%%%%%%%%%%%%%%%%%%%%%%
\subsection{Absolute Value Derivative}
\label{subsec:absolute-value-derivative}
%%%%%%%%%%%%%%%%%%%%%%%%%%%%%%%%%%%%%%%%%%%%%%%%%%%%%%%%%%%%

For

\[
x\neq0,
\]

\[
\boxed{
\frac{d}{dx}|x|
=
\frac{x}{|x|}
=
\operatorname{sgn}(x).
}
\]

At \(x=0\),

\[
\boxed{
|x|
}
\]

is not differentiable.

%%%%%%%%%%%%%%%%%%%%%%%%%%%%%%%%%%%%%%%%%%%%%%%%%%%%%%%%%%%%
\subsection{Sign Function}
\label{subsec:sign-function-derivative-reference}
%%%%%%%%%%%%%%%%%%%%%%%%%%%%%%%%%%%%%%%%%%%%%%%%%%%%%%%%%%%%

For \(x\neq0\),

\[
\boxed{
\operatorname{sgn}(x)
=
\begin{cases}
1, & x>0,\\
-1, & x<0.
\end{cases}
}
\]

Thus

\[
\boxed{
\frac{d}{dx}|x|
=
\operatorname{sgn}(x)
}
\]

away from zero.

%%%%%%%%%%%%%%%%%%%%%%%%%%%%%%%%%%%%%%%%%%%%%%%%%%%%%%%%%%%%
\subsection{Logarithmic Differentiation}
\label{subsec:logarithmic-differentiation-reference}
%%%%%%%%%%%%%%%%%%%%%%%%%%%%%%%%%%%%%%%%%%%%%%%%%%%%%%%%%%%%

Logarithmic differentiation is useful when variables occur in both base and
exponent or when many products and powers are present.

If

\[
y=f(x)>0,
\]

take logarithms:

\[
\ln y
=
\ln f(x).
\]

Differentiate implicitly:

\[
\boxed{
\frac{y'}{y}
=
\frac{d}{dx}
\ln f(x).
}
\]

Then solve for \(y'\).

%%%%%%%%%%%%%%%%%%%%%%%%%%%%%%%%%%%%%%%%%%%%%%%%%%%%%%%%%%%%
\subsection{Derivative of \(x^x\)}
\label{subsec:x-to-x-derivative}
%%%%%%%%%%%%%%%%%%%%%%%%%%%%%%%%%%%%%%%%%%%%%%%%%%%%%%%%%%%%

Let

\[
y=x^x,
\qquad
x>0.
\]

Take logarithms:

\[
\ln y
=
x\ln x.
\]

Differentiate:

\[
\frac{y'}{y}
=
\ln x+1.
\]

Therefore,

\[
\boxed{
\frac{d}{dx}x^x
=
x^x(\ln x+1).
}
\]

%%%%%%%%%%%%%%%%%%%%%%%%%%%%%%%%%%%%%%%%%%%%%%%%%%%%%%%%%%%%
\subsection{General Variable Power}
\label{subsec:general-variable-power-derivative}
%%%%%%%%%%%%%%%%%%%%%%%%%%%%%%%%%%%%%%%%%%%%%%%%%%%%%%%%%%%%

For positive \(f(x)\),

\[
y
=
[f(x)]^{g(x)}.
\]

Then

\[
\ln y
=
g(x)\ln f(x).
\]

Therefore,

\[
\boxed{
y'
=
f(x)^{g(x)}
\left[
g'(x)\ln f(x)
+
g(x)\frac{f'(x)}{f(x)}
\right].
}
\]

%%%%%%%%%%%%%%%%%%%%%%%%%%%%%%%%%%%%%%%%%%%%%%%%%%%%%%%%%%%%
\subsection{Implicit Differentiation}
\label{subsec:implicit-differentiation-reference}
%%%%%%%%%%%%%%%%%%%%%%%%%%%%%%%%%%%%%%%%%%%%%%%%%%%%%%%%%%%%

If \(x\) and \(y\) satisfy

\[
F(x,y)=0,
\]

differentiate both sides with respect to \(x\).

Treat

\[
y=y(x).
\]

For example,

\[
x^2+y^2=1
\]

gives

\[
2x
+
2y\frac{dy}{dx}
=
0.
\]

Thus

\[
\boxed{
\frac{dy}{dx}
=
-\frac{x}{y}.
}
\]

%%%%%%%%%%%%%%%%%%%%%%%%%%%%%%%%%%%%%%%%%%%%%%%%%%%%%%%%%%%%
\subsection{Implicit Function Formula}
\label{subsec:implicit-function-derivative-formula}
%%%%%%%%%%%%%%%%%%%%%%%%%%%%%%%%%%%%%%%%%%%%%%%%%%%%%%%%%%%%

If

\[
F(x,y)=0
\]

and

\[
F_y\neq0,
\]

then locally,

\[
\boxed{
\frac{dy}{dx}
=
-\frac{
F_x
}{
F_y
}.
}
\]

%%%%%%%%%%%%%%%%%%%%%%%%%%%%%%%%%%%%%%%%%%%%%%%%%%%%%%%%%%%%
\subsection{Higher Derivatives}
\label{subsec:higher-derivatives-reference}
%%%%%%%%%%%%%%%%%%%%%%%%%%%%%%%%%%%%%%%%%%%%%%%%%%%%%%%%%%%%

The second derivative is

\[
\boxed{
f''(x)
=
\frac{d^2f}{dx^2}.
}
\]

The third derivative is

\[
\boxed{
f'''(x).
}
\]

The \(n\)-th derivative is

\[
\boxed{
f^{(n)}(x).
}
\]

%%%%%%%%%%%%%%%%%%%%%%%%%%%%%%%%%%%%%%%%%%%%%%%%%%%%%%%%%%%%
\subsection{Worked Example: Higher Derivatives}
\label{subsec:worked-higher-derivatives}
%%%%%%%%%%%%%%%%%%%%%%%%%%%%%%%%%%%%%%%%%%%%%%%%%%%%%%%%%%%%

\begin{workedexamplebox}{Differentiate Repeatedly}

Let

\[
f(x)=x^4.
\]

Then

\[
f'(x)=4x^3,
\]

\[
f''(x)=12x^2,
\]

\[
f'''(x)=24x,
\]

and

\[
f^{(4)}(x)=24.
\]

Thus

\begin{answerbox}
\[
\boxed{
f^{(5)}(x)=0.
}
\]
\end{answerbox}

\end{workedexamplebox}

%%%%%%%%%%%%%%%%%%%%%%%%%%%%%%%%%%%%%%%%%%%%%%%%%%%%%%%%%%%%
\subsection{Derivative of an Inverse Function}
\label{subsec:inverse-function-derivative-general}
%%%%%%%%%%%%%%%%%%%%%%%%%%%%%%%%%%%%%%%%%%%%%%%%%%%%%%%%%%%%

If \(f\) is differentiable and invertible and

\[
f'(f^{-1}(y))\neq0,
\]

then

\[
\boxed{
(f^{-1})'(y)
=
\frac{
1
}{
f'(f^{-1}(y))
}.
}
\]

Equivalently,

\[
\boxed{
(f^{-1})'(x)
=
\frac{
1
}{
f'(f^{-1}(x))
}.
}
\]

%%%%%%%%%%%%%%%%%%%%%%%%%%%%%%%%%%%%%%%%%%%%%%%%%%%%%%%%%%%%
\subsection{Parametric First Derivative}
\label{subsec:parametric-first-derivative}
%%%%%%%%%%%%%%%%%%%%%%%%%%%%%%%%%%%%%%%%%%%%%%%%%%%%%%%%%%%%

If

\[
x=x(t),
\qquad
y=y(t),
\]

and

\[
\frac{dx}{dt}\neq0,
\]

then

\[
\boxed{
\frac{dy}{dx}
=
\frac{
dy/dt
}{
dx/dt
}.
}
\]

%%%%%%%%%%%%%%%%%%%%%%%%%%%%%%%%%%%%%%%%%%%%%%%%%%%%%%%%%%%%
\subsection{Parametric Second Derivative}
\label{subsec:parametric-second-derivative}
%%%%%%%%%%%%%%%%%%%%%%%%%%%%%%%%%%%%%%%%%%%%%%%%%%%%%%%%%%%%

\[
\boxed{
\frac{d^2y}{dx^2}
=
\frac{
\dfrac{d}{dt}
\left(
\dfrac{dy}{dx}
\right)
}{
dx/dt
}.
}
\]

%%%%%%%%%%%%%%%%%%%%%%%%%%%%%%%%%%%%%%%%%%%%%%%%%%%%%%%%%%%%
\subsection{Worked Example: Parametric Derivative}
\label{subsec:worked-parametric-derivative}
%%%%%%%%%%%%%%%%%%%%%%%%%%%%%%%%%%%%%%%%%%%%%%%%%%%%%%%%%%%%

\begin{workedexamplebox}{Differentiate Parametric Equations}

Let

\[
x=t^2,
\qquad
y=t^3.
\]

Then

\[
\frac{dx}{dt}=2t,
\]

and

\[
\frac{dy}{dt}=3t^2.
\]

Therefore, for \(t\neq0\),

\begin{answerbox}
\[
\boxed{
\frac{dy}{dx}
=
\frac{3t}{2}.
}
\]
\end{answerbox}

\end{workedexamplebox}

%%%%%%%%%%%%%%%%%%%%%%%%%%%%%%%%%%%%%%%%%%%%%%%%%%%%%%%%%%%%
\subsection{Polar First Derivative}
\label{subsec:polar-first-derivative}
%%%%%%%%%%%%%%%%%%%%%%%%%%%%%%%%%%%%%%%%%%%%%%%%%%%%%%%%%%%%

For

\[
r=r(\theta),
\]

with

\[
x=r\cos\theta,
\qquad
y=r\sin\theta,
\]

we have

\[
\boxed{
\frac{dy}{dx}
=
\frac{
r'(\theta)\sin\theta
+
r(\theta)\cos\theta
}{
r'(\theta)\cos\theta
-
r(\theta)\sin\theta
}.
}
\]

%%%%%%%%%%%%%%%%%%%%%%%%%%%%%%%%%%%%%%%%%%%%%%%%%%%%%%%%%%%%
\subsection{Vector-Valued Derivatives}
\label{subsec:vector-valued-derivative}
%%%%%%%%%%%%%%%%%%%%%%%%%%%%%%%%%%%%%%%%%%%%%%%%%%%%%%%%%%%%

For

\[
\mathbf{r}(t)
=
\langle
x(t),
y(t),
z(t)
\rangle,
\]

\[
\boxed{
\mathbf{r}'(t)
=
\langle
x'(t),
y'(t),
z'(t)
\rangle.
}
\]

The derivative is taken componentwise.

%%%%%%%%%%%%%%%%%%%%%%%%%%%%%%%%%%%%%%%%%%%%%%%%%%%%%%%%%%%%
\subsection{Velocity and Acceleration}
\label{subsec:velocity-acceleration-derivative}
%%%%%%%%%%%%%%%%%%%%%%%%%%%%%%%%%%%%%%%%%%%%%%%%%%%%%%%%%%%%

If

\[
\mathbf{r}(t)
\]

is position, then velocity is

\[
\boxed{
\mathbf{v}(t)
=
\mathbf{r}'(t).
}
\]

Acceleration is

\[
\boxed{
\mathbf{a}(t)
=
\mathbf{r}''(t).
}
\]

Speed is

\[
\boxed{
\|\mathbf{v}(t)\|.
}
\]

%%%%%%%%%%%%%%%%%%%%%%%%%%%%%%%%%%%%%%%%%%%%%%%%%%%%%%%%%%%%
\subsection{Partial Derivative}
\label{subsec:partial-derivative-table-reference}
%%%%%%%%%%%%%%%%%%%%%%%%%%%%%%%%%%%%%%%%%%%%%%%%%%%%%%%%%%%%

For

\[
f=f(x_1,\ldots,x_n),
\]

the partial derivative with respect to \(x_i\) is

\[
\boxed{
\frac{\partial f}{\partial x_i}.
}
\]

Other variables are held fixed.

%%%%%%%%%%%%%%%%%%%%%%%%%%%%%%%%%%%%%%%%%%%%%%%%%%%%%%%%%%%%
\subsection{Gradient}
\label{subsec:gradient-derivative-table}
%%%%%%%%%%%%%%%%%%%%%%%%%%%%%%%%%%%%%%%%%%%%%%%%%%%%%%%%%%%%

For scalar function

\[
f:\mathbb{R}^n\to\mathbb{R},
\]

the gradient is

\[
\boxed{
\nabla f
=
\begin{bmatrix}
\frac{\partial f}{\partial x_1}\\
\vdots\\
\frac{\partial f}{\partial x_n}
\end{bmatrix}.
}
\]

%%%%%%%%%%%%%%%%%%%%%%%%%%%%%%%%%%%%%%%%%%%%%%%%%%%%%%%%%%%%
\subsection{Directional Derivative}
\label{subsec:directional-derivative-reference}
%%%%%%%%%%%%%%%%%%%%%%%%%%%%%%%%%%%%%%%%%%%%%%%%%%%%%%%%%%%%

For unit vector \(\mathbf{u}\),

\[
\boxed{
D_{\mathbf{u}}f
=
\nabla f\cdot\mathbf{u}.
}
\]

This gives the instantaneous rate of change of \(f\) in direction
\(\mathbf{u}\).

%%%%%%%%%%%%%%%%%%%%%%%%%%%%%%%%%%%%%%%%%%%%%%%%%%%%%%%%%%%%
\subsection{Jacobian Matrix}
\label{subsec:jacobian-derivative-reference}
%%%%%%%%%%%%%%%%%%%%%%%%%%%%%%%%%%%%%%%%%%%%%%%%%%%%%%%%%%%%

For

\[
F:\mathbb{R}^n\to\mathbb{R}^m,
\]

with components

\[
F_1,\ldots,F_m,
\]

the Jacobian is

\[
\boxed{
J_F
=
\begin{bmatrix}
\dfrac{\partial F_1}{\partial x_1}
&
\cdots
&
\dfrac{\partial F_1}{\partial x_n}
\\
\vdots
&
\ddots
&
\vdots
\\
\dfrac{\partial F_m}{\partial x_1}
&
\cdots
&
\dfrac{\partial F_m}{\partial x_n}
\end{bmatrix}.
}
\]

%%%%%%%%%%%%%%%%%%%%%%%%%%%%%%%%%%%%%%%%%%%%%%%%%%%%%%%%%%%%
\subsection{Hessian Matrix}
\label{subsec:hessian-derivative-reference}
%%%%%%%%%%%%%%%%%%%%%%%%%%%%%%%%%%%%%%%%%%%%%%%%%%%%%%%%%%%%

For twice differentiable scalar function

\[
f:\mathbb{R}^n\to\mathbb{R},
\]

the Hessian is

\[
\boxed{
H_f
=
\nabla^2f
=
\left[
\frac{
\partial^2f
}{
\partial x_i\partial x_j
}
\right].
}
\]

%%%%%%%%%%%%%%%%%%%%%%%%%%%%%%%%%%%%%%%%%%%%%%%%%%%%%%%%%%%%
\subsection{Mixed Partial Derivatives}
\label{subsec:mixed-partial-reference}
%%%%%%%%%%%%%%%%%%%%%%%%%%%%%%%%%%%%%%%%%%%%%%%%%%%%%%%%%%%%

Under appropriate regularity conditions,

\[
\boxed{
\frac{
\partial^2f
}{
\partial x\partial y
}
=
\frac{
\partial^2f
}{
\partial y\partial x
}.
}
\]

This is often called Clairaut's or Schwarz's theorem.

%%%%%%%%%%%%%%%%%%%%%%%%%%%%%%%%%%%%%%%%%%%%%%%%%%%%%%%%%%%%
\subsection{Multivariable Chain Rule}
\label{subsec:multivariable-chain-rule-reference}
%%%%%%%%%%%%%%%%%%%%%%%%%%%%%%%%%%%%%%%%%%%%%%%%%%%%%%%%%%%%

If

\[
z=f(x,y),
\]

where

\[
x=x(t),
\qquad
y=y(t),
\]

then

\[
\boxed{
\frac{dz}{dt}
=
\frac{\partial f}{\partial x}
\frac{dx}{dt}
+
\frac{\partial f}{\partial y}
\frac{dy}{dt}.
}
\]

%%%%%%%%%%%%%%%%%%%%%%%%%%%%%%%%%%%%%%%%%%%%%%%%%%%%%%%%%%%%
\subsection{Matrix Form of the Chain Rule}
\label{subsec:matrix-chain-rule-reference}
%%%%%%%%%%%%%%%%%%%%%%%%%%%%%%%%%%%%%%%%%%%%%%%%%%%%%%%%%%%%

For differentiable maps

\[
F:\mathbb{R}^n\to\mathbb{R}^m
\]

and

\[
G:\mathbb{R}^p\to\mathbb{R}^n,
\]

\[
\boxed{
J_{F\circ G}(x)
=
J_F(G(x))
J_G(x).
}
\]

%%%%%%%%%%%%%%%%%%%%%%%%%%%%%%%%%%%%%%%%%%%%%%%%%%%%%%%%%%%%
\subsection{Derivative of a Dot Product}
\label{subsec:dot-product-derivative}
%%%%%%%%%%%%%%%%%%%%%%%%%%%%%%%%%%%%%%%%%%%%%%%%%%%%%%%%%%%%

For differentiable vector functions

\[
\mathbf{u}(t),
\qquad
\mathbf{v}(t),
\]

\[
\boxed{
\frac{d}{dt}
(\mathbf{u}\cdot\mathbf{v})
=
\mathbf{u}'\cdot\mathbf{v}
+
\mathbf{u}\cdot\mathbf{v}'.
}
\]

%%%%%%%%%%%%%%%%%%%%%%%%%%%%%%%%%%%%%%%%%%%%%%%%%%%%%%%%%%%%
\subsection{Derivative of a Cross Product}
\label{subsec:cross-product-derivative}
%%%%%%%%%%%%%%%%%%%%%%%%%%%%%%%%%%%%%%%%%%%%%%%%%%%%%%%%%%%%

\[
\boxed{
\frac{d}{dt}
(\mathbf{u}\times\mathbf{v})
=
\mathbf{u}'\times\mathbf{v}
+
\mathbf{u}\times\mathbf{v}'.
}
\]

%%%%%%%%%%%%%%%%%%%%%%%%%%%%%%%%%%%%%%%%%%%%%%%%%%%%%%%%%%%%
\subsection{Derivative of a Vector Norm Squared}
\label{subsec:norm-squared-derivative}
%%%%%%%%%%%%%%%%%%%%%%%%%%%%%%%%%%%%%%%%%%%%%%%%%%%%%%%%%%%%

Since

\[
\|\mathbf{x}(t)\|^2
=
\mathbf{x}\cdot\mathbf{x},
\]

\[
\boxed{
\frac{d}{dt}
\|\mathbf{x}\|^2
=
2\mathbf{x}\cdot\mathbf{x}'.
}
\]

%%%%%%%%%%%%%%%%%%%%%%%%%%%%%%%%%%%%%%%%%%%%%%%%%%%%%%%%%%%%
\subsection{Derivative of a Vector Norm}
\label{subsec:norm-derivative}
%%%%%%%%%%%%%%%%%%%%%%%%%%%%%%%%%%%%%%%%%%%%%%%%%%%%%%%%%%%%

If

\[
\mathbf{x}(t)\neq\mathbf{0},
\]

then

\[
\boxed{
\frac{d}{dt}
\|\mathbf{x}\|
=
\frac{
\mathbf{x}\cdot\mathbf{x}'
}{
\|\mathbf{x}\|
}.
}
\]

%%%%%%%%%%%%%%%%%%%%%%%%%%%%%%%%%%%%%%%%%%%%%%%%%%%%%%%%%%%%
\subsection{Derivative of Matrix Inverse}
\label{subsec:matrix-inverse-derivative}
%%%%%%%%%%%%%%%%%%%%%%%%%%%%%%%%%%%%%%%%%%%%%%%%%%%%%%%%%%%%

If \(A(t)\) is differentiable and invertible,

\[
\boxed{
\frac{d}{dt}
A^{-1}
=
-A^{-1}A'A^{-1}.
}
\]

This follows from differentiating

\[
AA^{-1}=I.
\]

%%%%%%%%%%%%%%%%%%%%%%%%%%%%%%%%%%%%%%%%%%%%%%%%%%%%%%%%%%%%
\subsection{Derivative of a Determinant}
\label{subsec:determinant-derivative}
%%%%%%%%%%%%%%%%%%%%%%%%%%%%%%%%%%%%%%%%%%%%%%%%%%%%%%%%%%%%

For differentiable invertible matrix \(A(t)\),

\[
\boxed{
\frac{d}{dt}
\det A
=
\det(A)
\operatorname{tr}
\left(
A^{-1}A'
\right).
}
\]

This is Jacobi's formula.

%%%%%%%%%%%%%%%%%%%%%%%%%%%%%%%%%%%%%%%%%%%%%%%%%%%%%%%%%%%%
\subsection{Derivative of Log Determinant}
\label{subsec:logdet-derivative}
%%%%%%%%%%%%%%%%%%%%%%%%%%%%%%%%%%%%%%%%%%%%%%%%%%%%%%%%%%%%

For differentiable invertible \(A(t)\) with positive determinant,

\[
\boxed{
\frac{d}{dt}
\ln\det A
=
\operatorname{tr}
\left(
A^{-1}A'
\right).
}
\]

%%%%%%%%%%%%%%%%%%%%%%%%%%%%%%%%%%%%%%%%%%%%%%%%%%%%%%%%%%%%
\subsection{Gradient of a Linear Form}
\label{subsec:gradient-linear-form}
%%%%%%%%%%%%%%%%%%%%%%%%%%%%%%%%%%%%%%%%%%%%%%%%%%%%%%%%%%%%

For constant vector \(\mathbf{a}\),

\[
f(\mathbf{x})
=
\mathbf{a}^T\mathbf{x},
\]

then

\[
\boxed{
\nabla f
=
\mathbf{a}.
}
\]

%%%%%%%%%%%%%%%%%%%%%%%%%%%%%%%%%%%%%%%%%%%%%%%%%%%%%%%%%%%%
\subsection{Gradient of a Quadratic Norm}
\label{subsec:gradient-norm-squared}
%%%%%%%%%%%%%%%%%%%%%%%%%%%%%%%%%%%%%%%%%%%%%%%%%%%%%%%%%%%%

For

\[
f(\mathbf{x})
=
\mathbf{x}^T\mathbf{x},
\]

\[
\boxed{
\nabla f
=
2\mathbf{x}.
}
\]

%%%%%%%%%%%%%%%%%%%%%%%%%%%%%%%%%%%%%%%%%%%%%%%%%%%%%%%%%%%%
\subsection{Gradient of a Quadratic Form}
\label{subsec:gradient-quadratic-form}
%%%%%%%%%%%%%%%%%%%%%%%%%%%%%%%%%%%%%%%%%%%%%%%%%%%%%%%%%%%%

For

\[
f(\mathbf{x})
=
\mathbf{x}^TA\mathbf{x},
\]

\[
\boxed{
\nabla f
=
(A+A^T)\mathbf{x}.
}
\]

If \(A\) is symmetric,

\[
\boxed{
\nabla f
=
2A\mathbf{x}.
}
\]

%%%%%%%%%%%%%%%%%%%%%%%%%%%%%%%%%%%%%%%%%%%%%%%%%%%%%%%%%%%%
\subsection{Hessian of a Quadratic Form}
\label{subsec:hessian-quadratic-form}
%%%%%%%%%%%%%%%%%%%%%%%%%%%%%%%%%%%%%%%%%%%%%%%%%%%%%%%%%%%%

For

\[
f(\mathbf{x})
=
\mathbf{x}^TA\mathbf{x},
\]

\[
\boxed{
\nabla^2f
=
A+A^T.
}
\]

If \(A=A^T\),

\[
\boxed{
\nabla^2f
=
2A.
}
\]

%%%%%%%%%%%%%%%%%%%%%%%%%%%%%%%%%%%%%%%%%%%%%%%%%%%%%%%%%%%%
\subsection{Gradient of Least Squares}
\label{subsec:gradient-least-squares}
%%%%%%%%%%%%%%%%%%%%%%%%%%%%%%%%%%%%%%%%%%%%%%%%%%%%%%%%%%%%

Let

\[
f(\mathbf{x})
=
\frac12
\|A\mathbf{x}-\mathbf{b}\|^2.
\]

Then

\[
\boxed{
\nabla f
=
A^T
(A\mathbf{x}-\mathbf{b}).
}
\]

The Hessian is

\[
\boxed{
\nabla^2f
=
A^TA.
}
\]

%%%%%%%%%%%%%%%%%%%%%%%%%%%%%%%%%%%%%%%%%%%%%%%%%%%%%%%%%%%%
\subsection{Derivative of Logistic Function}
\label{subsec:logistic-derivative-reference}
%%%%%%%%%%%%%%%%%%%%%%%%%%%%%%%%%%%%%%%%%%%%%%%%%%%%%%%%%%%%

For

\[
\sigma(x)
=
\frac{1}{1+e^{-x}},
\]

\[
\boxed{
\sigma'(x)
=
\sigma(x)
[1-\sigma(x)].
}
\]

This derivative appears in statistics, dynamical systems, and machine
learning.

%%%%%%%%%%%%%%%%%%%%%%%%%%%%%%%%%%%%%%%%%%%%%%%%%%%%%%%%%%%%
\subsection{Derivative of Softplus}
\label{subsec:softplus-derivative}
%%%%%%%%%%%%%%%%%%%%%%%%%%%%%%%%%%%%%%%%%%%%%%%%%%%%%%%%%%%%

For

\[
f(x)
=
\ln(1+e^x),
\]

\[
\boxed{
f'(x)
=
\frac{1}{1+e^{-x}}.
}
\]

Thus the derivative of softplus is the logistic function.

%%%%%%%%%%%%%%%%%%%%%%%%%%%%%%%%%%%%%%%%%%%%%%%%%%%%%%%%%%%%
\subsection{Derivative of Gaussian Function}
\label{subsec:gaussian-function-derivative}
%%%%%%%%%%%%%%%%%%%%%%%%%%%%%%%%%%%%%%%%%%%%%%%%%%%%%%%%%%%%

For

\[
f(x)
=
e^{-x^2},
\]

\[
\boxed{
f'(x)
=
-2xe^{-x^2}.
}
\]

More generally,

\[
f(x)
=
e^{-ax^2}
\]

gives

\[
\boxed{
f'(x)
=
-2ax e^{-ax^2}.
}
\]

%%%%%%%%%%%%%%%%%%%%%%%%%%%%%%%%%%%%%%%%%%%%%%%%%%%%%%%%%%%%
\subsection{Derivative of General Gaussian Density}
\label{subsec:normal-density-derivative}
%%%%%%%%%%%%%%%%%%%%%%%%%%%%%%%%%%%%%%%%%%%%%%%%%%%%%%%%%%%%

For

\[
f(x)
=
\frac{
1
}{
\sigma\sqrt{2\pi}
}
\exp
\left(
-\frac{
(x-\mu)^2
}{
2\sigma^2
}
\right),
\]

\[
\boxed{
f'(x)
=
-\frac{
x-\mu
}{
\sigma^2
}
f(x).
}
\]

%%%%%%%%%%%%%%%%%%%%%%%%%%%%%%%%%%%%%%%%%%%%%%%%%%%%%%%%%%%%
\subsection{Derivative of Power-Exponential Product}
\label{subsec:power-exponential-product-derivative}
%%%%%%%%%%%%%%%%%%%%%%%%%%%%%%%%%%%%%%%%%%%%%%%%%%%%%%%%%%%%

For

\[
f(x)
=
x^ne^{ax},
\]

\[
\boxed{
f'(x)
=
nx^{n-1}e^{ax}
+
ax^ne^{ax}.
}
\]

Factoring,

\[
\boxed{
f'(x)
=
x^{n-1}e^{ax}
(n+ax).
}
\]

%%%%%%%%%%%%%%%%%%%%%%%%%%%%%%%%%%%%%%%%%%%%%%%%%%%%%%%%%%%%
\subsection{Derivative of Exponential Trigonometric Product}
\label{subsec:exp-trig-product-derivative}
%%%%%%%%%%%%%%%%%%%%%%%%%%%%%%%%%%%%%%%%%%%%%%%%%%%%%%%%%%%%

For

\[
f(x)
=
e^{ax}\sin bx,
\]

\[
\boxed{
f'(x)
=
e^{ax}
(a\sin bx+b\cos bx).
}
\]

For

\[
f(x)
=
e^{ax}\cos bx,
\]

\[
\boxed{
f'(x)
=
e^{ax}
(a\cos bx-b\sin bx).
}
\]

%%%%%%%%%%%%%%%%%%%%%%%%%%%%%%%%%%%%%%%%%%%%%%%%%%%%%%%%%%%%
\subsection{Derivative of \(x\ln x\)}
\label{subsec:x-log-x-derivative}
%%%%%%%%%%%%%%%%%%%%%%%%%%%%%%%%%%%%%%%%%%%%%%%%%%%%%%%%%%%%

For \(x>0\),

\[
\boxed{
\frac{d}{dx}
(x\ln x)
=
\ln x+1.
}
\]

%%%%%%%%%%%%%%%%%%%%%%%%%%%%%%%%%%%%%%%%%%%%%%%%%%%%%%%%%%%%
\subsection{Derivative of \(\ln(x^2+a^2)\)}
\label{subsec:log-square-plus-constant-derivative}
%%%%%%%%%%%%%%%%%%%%%%%%%%%%%%%%%%%%%%%%%%%%%%%%%%%%%%%%%%%%

For constant \(a\),

\[
\boxed{
\frac{d}{dx}
\ln(x^2+a^2)
=
\frac{
2x
}{
x^2+a^2
}.
}
\]

%%%%%%%%%%%%%%%%%%%%%%%%%%%%%%%%%%%%%%%%%%%%%%%%%%%%%%%%%%%%
\subsection{Derivative of \(\sqrt{a^2-x^2}\)}
\label{subsec:sqrt-a2-minus-x2-derivative}
%%%%%%%%%%%%%%%%%%%%%%%%%%%%%%%%%%%%%%%%%%%%%%%%%%%%%%%%%%%%

Where defined,

\[
\boxed{
\frac{d}{dx}
\sqrt{a^2-x^2}
=
-\frac{
x
}{
\sqrt{a^2-x^2}
}.
}
\]

%%%%%%%%%%%%%%%%%%%%%%%%%%%%%%%%%%%%%%%%%%%%%%%%%%%%%%%%%%%%
\subsection{Derivative of \(\sqrt{x^2+a^2}\)}
\label{subsec:sqrt-x2-plus-a2-derivative}
%%%%%%%%%%%%%%%%%%%%%%%%%%%%%%%%%%%%%%%%%%%%%%%%%%%%%%%%%%%%

\[
\boxed{
\frac{d}{dx}
\sqrt{x^2+a^2}
=
\frac{
x
}{
\sqrt{x^2+a^2}
}.
}
\]

%%%%%%%%%%%%%%%%%%%%%%%%%%%%%%%%%%%%%%%%%%%%%%%%%%%%%%%%%%%%
\subsection{Derivative of \(\ln|\sec x+\tan x|\)}
\label{subsec:log-sec-tan-derivative}
%%%%%%%%%%%%%%%%%%%%%%%%%%%%%%%%%%%%%%%%%%%%%%%%%%%%%%%%%%%%

Where defined,

\[
\boxed{
\frac{d}{dx}
\ln|\sec x+\tan x|
=
\sec x.
}
\]

This identity is useful when deriving the antiderivative of secant.

%%%%%%%%%%%%%%%%%%%%%%%%%%%%%%%%%%%%%%%%%%%%%%%%%%%%%%%%%%%%
\subsection{Derivative of \(\ln|\csc x-\cot x|\)}
\label{subsec:log-csc-cot-derivative}
%%%%%%%%%%%%%%%%%%%%%%%%%%%%%%%%%%%%%%%%%%%%%%%%%%%%%%%%%%%%

Where defined,

\[
\boxed{
\frac{d}{dx}
\ln|\csc x-\cot x|
=
\csc x.
}
\]

%%%%%%%%%%%%%%%%%%%%%%%%%%%%%%%%%%%%%%%%%%%%%%%%%%%%%%%%%%%%
\subsection{Derivative of a General Composite}
\label{subsec:general-composite-derivative}
%%%%%%%%%%%%%%%%%%%%%%%%%%%%%%%%%%%%%%%%%%%%%%%%%%%%%%%%%%%%

For

\[
y
=
F(G(H(x))),
\]

\[
\boxed{
y'
=
F'(G(H(x)))
G'(H(x))
H'(x).
}
\]

Each nested layer contributes one chain-rule factor.

%%%%%%%%%%%%%%%%%%%%%%%%%%%%%%%%%%%%%%%%%%%%%%%%%%%%%%%%%%%%
\subsection{Worked Example: Nested Composition}
\label{subsec:worked-nested-composition}
%%%%%%%%%%%%%%%%%%%%%%%%%%%%%%%%%%%%%%%%%%%%%%%%%%%%%%%%%%%%

\begin{workedexamplebox}{Differentiate Several Layers}

Differentiate

\[
y
=
e^{\sin(x^2)}.
\]

Differentiate the outer exponential:

\[
e^{\sin(x^2)}.
\]

Differentiate the sine:

\[
\cos(x^2).
\]

Differentiate \(x^2\):

\[
2x.
\]

Therefore,

\begin{answerbox}
\[
\boxed{
y'
=
2x
e^{\sin(x^2)}
\cos(x^2).
}
\]
\end{answerbox}

\end{workedexamplebox}

%%%%%%%%%%%%%%%%%%%%%%%%%%%%%%%%%%%%%%%%%%%%%%%%%%%%%%%%%%%%
\subsection{Quick Core Derivative Table}
\label{subsec:quick-core-derivative-table}
%%%%%%%%%%%%%%%%%%%%%%%%%%%%%%%%%%%%%%%%%%%%%%%%%%%%%%%%%%%%

\[
\boxed{
(c)'=0
}
\]

\[
\boxed{
(x^n)'=nx^{n-1}
}
\]

\[
\boxed{
(e^x)'=e^x
}
\]

\[
\boxed{
(a^x)'=a^x\ln a
}
\]

\[
\boxed{
(\ln x)'=\frac1x
}
\]

\[
\boxed{
(\log_a x)'
=
\frac{1}{x\ln a}
}
\]

\[
\boxed{
(\sqrt{x})'
=
\frac{1}{2\sqrt{x}}
}
\]

%%%%%%%%%%%%%%%%%%%%%%%%%%%%%%%%%%%%%%%%%%%%%%%%%%%%%%%%%%%%
\subsection{Quick Rule Table}
\label{subsec:quick-derivative-rule-table}
%%%%%%%%%%%%%%%%%%%%%%%%%%%%%%%%%%%%%%%%%%%%%%%%%%%%%%%%%%%%

\[
\boxed{
(cf)'=cf'
}
\]

\[
\boxed{
(f+g)'=f'+g'
}
\]

\[
\boxed{
(f-g)'=f'-g'
}
\]

\[
\boxed{
(fg)'=f'g+fg'
}
\]

\[
\boxed{
\left(
\frac fg
\right)'
=
\frac{
f'g-fg'
}{
g^2
}
}
\]

\[
\boxed{
(f\circ g)'
=
(f'\circ g)g'
}
\]

%%%%%%%%%%%%%%%%%%%%%%%%%%%%%%%%%%%%%%%%%%%%%%%%%%%%%%%%%%%%
\subsection{Common Errors}
\label{subsec:derivative-tables-common-errors}
%%%%%%%%%%%%%%%%%%%%%%%%%%%%%%%%%%%%%%%%%%%%%%%%%%%%%%%%%%%%

\subsubsection{Forgetting the Chain Rule}

In general,

\[
\frac{d}{dx}
\sin(x^2)
\]

is not simply

\[
\cos(x^2).
\]

The correct derivative is

\[
\boxed{
2x\cos(x^2).
}
\]

%%%%%%%%%%%%%%%%%%%%%%%%%%%%%%%%%%%%%%%%%%%%%%%%%%%%%%%%%%%%
\subsubsection{Incorrect Product Rule}

In general,

\[
\boxed{
(fg)'
\neq
f'g'.
}
\]

The correct formula is

\[
\boxed{
(fg)'
=
f'g+fg'.
}
\]

%%%%%%%%%%%%%%%%%%%%%%%%%%%%%%%%%%%%%%%%%%%%%%%%%%%%%%%%%%%%
\subsubsection{Incorrect Quotient Rule}

In general,

\[
\boxed{
\left(
\frac fg
\right)'
\neq
\frac{f'}{g'}.
}
\]

Use

\[
\boxed{
\frac{
f'g-fg'
}{
g^2
}.
}
\]

%%%%%%%%%%%%%%%%%%%%%%%%%%%%%%%%%%%%%%%%%%%%%%%%%%%%%%%%%%%%
\subsubsection{Forgetting the Negative Sign for Cosine}

\[
\boxed{
(\cos x)'
=
-\sin x.
}
\]

%%%%%%%%%%%%%%%%%%%%%%%%%%%%%%%%%%%%%%%%%%%%%%%%%%%%%%%%%%%%
\subsubsection{Forgetting the Negative Sign for Cotangent}

\[
\boxed{
(\cot x)'
=
-\csc^2x.
}
\]

%%%%%%%%%%%%%%%%%%%%%%%%%%%%%%%%%%%%%%%%%%%%%%%%%%%%%%%%%%%%
\subsubsection{Forgetting the Negative Sign for Cosecant}

\[
\boxed{
(\csc x)'
=
-\csc x\cot x.
}
\]

%%%%%%%%%%%%%%%%%%%%%%%%%%%%%%%%%%%%%%%%%%%%%%%%%%%%%%%%%%%%
\subsubsection{Confusing Inverse Trig With Reciprocal}

\[
\boxed{
\sin^{-1}x
=
\arcsin x
}
\]

in inverse-function notation.

It is not

\[
\frac1{\sin x}.
\]

The reciprocal is

\[
\boxed{
\csc x.
}
\]

%%%%%%%%%%%%%%%%%%%%%%%%%%%%%%%%%%%%%%%%%%%%%%%%%%%%%%%%%%%%
\subsubsection{Missing Absolute Values in Arcsec and Arccsc Derivatives}

The standard real formulas contain

\[
\boxed{
|x|.
}
\]

Thus

\[
\boxed{
(\operatorname{arcsec}x)'
=
\frac1{|x|\sqrt{x^2-1}}.
}
\]

%%%%%%%%%%%%%%%%%%%%%%%%%%%%%%%%%%%%%%%%%%%%%%%%%%%%%%%%%%%%
\subsubsection{Forgetting Absolute Value in \(\ln|u|\)}

The identity

\[
\boxed{
\frac{d}{dx}
\ln|u|
=
\frac{u'}{u}
}
\]

holds wherever \(u\neq0\).

Using only \(\ln u\) restricts to \(u>0\).

%%%%%%%%%%%%%%%%%%%%%%%%%%%%%%%%%%%%%%%%%%%%%%%%%%%%%%%%%%%%
\subsubsection{Differentiating \(x^x\) With the Ordinary Power Rule}

The power rule alone does not apply because both base and exponent vary.

Use logarithmic differentiation:

\[
\boxed{
(x^x)'
=
x^x(\ln x+1).
}
\]

%%%%%%%%%%%%%%%%%%%%%%%%%%%%%%%%%%%%%%%%%%%%%%%%%%%%%%%%%%%%
\subsubsection{Dropping \(dy/dx\) in Implicit Differentiation}

If \(y=y(x)\),

\[
\frac{d}{dx}y^2
\]

is

\[
\boxed{
2y\frac{dy}{dx},
}
\]

not simply \(2y\).

%%%%%%%%%%%%%%%%%%%%%%%%%%%%%%%%%%%%%%%%%%%%%%%%%%%%%%%%%%%%
\subsubsection{Assuming Differentiability From Continuity}

Differentiability implies continuity, but continuity does not imply
differentiability.

For example,

\[
\boxed{
|x|
}
\]

is continuous at \(0\) but not differentiable there.

%%%%%%%%%%%%%%%%%%%%%%%%%%%%%%%%%%%%%%%%%%%%%%%%%%%%%%%%%%%%
\subsubsection{Confusing Partial and Total Derivatives}

For \(f(x,y)\),

\[
\frac{\partial f}{\partial x}
\]

holds \(y\) fixed.

If \(y=y(x)\), then the total derivative involves chain-rule terms.

%%%%%%%%%%%%%%%%%%%%%%%%%%%%%%%%%%%%%%%%%%%%%%%%%%%%%%%%%%%%
\subsubsection{Ignoring Matrix Multiplication Order}

For matrix-valued expressions, multiplication may be noncommutative.

For example,

\[
\boxed{
\frac{d}{dt}
A^{-1}
=
-A^{-1}A'A^{-1}
}
\]

cannot generally be rearranged.

%%%%%%%%%%%%%%%%%%%%%%%%%%%%%%%%%%%%%%%%%%%%%%%%%%%%%%%%%%%%
\subsection{Derivative Evaluation Workflow}
\label{subsec:derivative-evaluation-workflow}
%%%%%%%%%%%%%%%%%%%%%%%%%%%%%%%%%%%%%%%%%%%%%%%%%%%%%%%%%%%%

A practical differentiation strategy is:

\begin{enumerate}
    \item identify the outermost operation;
    \item simplify algebra first when useful;
    \item classify the structure as sum, product, quotient, or composition;
    \item apply the appropriate rule;
    \item use the chain rule for every nontrivial inner function;
    \item preserve constants;
    \item simplify only after differentiating correctly;
    \item check signs in trigonometric derivatives;
    \item check domains for logarithms, radicals, and inverse functions;
    \item use logarithmic differentiation for variable exponents;
    \item use implicit differentiation when \(y\) is not isolated;
    \item verify dimensions or units when differentiating physical quantities.
\end{enumerate}

%%%%%%%%%%%%%%%%%%%%%%%%%%%%%%%%%%%%%%%%%%%%%%%%%%%%%%%%%%%%
\subsection{Exercises}
\label{subsec:derivative-table-exercises}
%%%%%%%%%%%%%%%%%%%%%%%%%%%%%%%%%%%%%%%%%%%%%%%%%%%%%%%%%%%%

\begin{exercise}[1]
Differentiate

\[
f(x)=7.
\]
\end{exercise}

\begin{exercise}[1]
Differentiate

\[
f(x)=x^8.
\]
\end{exercise}

\begin{exercise}[1]
Differentiate

\[
f(x)=3x^5-2x^2+7.
\]
\end{exercise}

\begin{exercise}[1]
Differentiate

\[
f(x)=\frac1x.
\]
\end{exercise}

\begin{exercise}[1]
Differentiate

\[
f(x)=\sqrt{x}.
\]
\end{exercise}

\begin{exercise}[2]
Differentiate

\[
f(x)
=
(x^2+3)^5.
\]
\end{exercise}

\begin{exercise}[2]
Differentiate

\[
f(x)
=
\left(
(x+1)^2+3
\right)^4.
\]
\end{exercise}

\begin{exercise}[2]
Differentiate

\[
f(x)
=
(x^3+1)(x^2-4).
\]
\end{exercise}

\begin{exercise}[2]
Differentiate

\[
f(x)
=
\frac{x^2+1}{x-2}.
\]
\end{exercise}

\begin{exercise}[2]
Differentiate

\[
f(x)=e^{3x}.
\]
\end{exercise}

\begin{exercise}[2]
Differentiate

\[
f(x)=2^x.
\]
\end{exercise}

\begin{exercise}[2]
Differentiate

\[
f(x)=\ln(x^2+1).
\]
\end{exercise}

\begin{exercise}[2]
Differentiate

\[
f(x)=\log_5(3x+1).
\]
\end{exercise}

\begin{exercise}[2]
Differentiate

\[
f(x)=\sin(4x).
\]
\end{exercise}

\begin{exercise}[2]
Differentiate

\[
f(x)=\cos(x^2).
\]
\end{exercise}

\begin{exercise}[2]
Differentiate

\[
f(x)=\tan(5x-1).
\]
\end{exercise}

\begin{exercise}[2]
Differentiate

\[
f(x)=\sec(x^3).
\]
\end{exercise}

\begin{exercise}[2]
Differentiate

\[
f(x)=\csc(2x).
\]
\end{exercise}

\begin{exercise}[2]
Differentiate

\[
f(x)=\cot(7x).
\]
\end{exercise}

\begin{exercise}[2]
Differentiate

\[
f(x)=\arcsin(2x).
\]
\end{exercise}

\begin{exercise}[2]
Differentiate

\[
f(x)=\arccos(x^2).
\]
\end{exercise}

\begin{exercise}[2]
Differentiate

\[
f(x)=\arctan(3x).
\]
\end{exercise}

\begin{exercise}[2]
Differentiate

\[
f(x)=\sinh(4x).
\]
\end{exercise}

\begin{exercise}[2]
Differentiate

\[
f(x)=\cosh(x^2).
\]
\end{exercise}

\begin{exercise}[2]
Differentiate

\[
f(x)=\tanh(5x).
\]
\end{exercise}

\begin{exercise}[2]
Differentiate

\[
f(x)=e^{\sin x}.
\]
\end{exercise}

\begin{exercise}[2]
Differentiate

\[
f(x)=\sin(e^x).
\]
\end{exercise}

\begin{exercise}[2]
Differentiate

\[
f(x)=\ln(\sin x).
\]
\end{exercise}

\begin{exercise}[2]
Differentiate

\[
f(x)=\sqrt{1-x^2}.
\]
\end{exercise}

\begin{exercise}[2]
Differentiate

\[
f(x)=x\ln x.
\]
\end{exercise}

\begin{exercise}[2]
Differentiate

\[
f(x)=x^x.
\]
\end{exercise}

\begin{exercise}[3]
Use logarithmic differentiation to differentiate

\[
y
=
\frac{
(x^2+1)^3
\sqrt{x-1}
}{
(x+2)^5
}.
\]
\end{exercise}

\begin{exercise}[3]
Differentiate implicitly:

\[
x^2+y^2=25.
\]
\end{exercise}

\begin{exercise}[3]
Differentiate implicitly:

\[
x^3+y^3=6xy.
\]
\end{exercise}

\begin{exercise}[3]
Use

\[
F_x/F_y
\]

to find \(dy/dx\) from an implicitly defined curve.
\end{exercise}

\begin{exercise}[3]
Find the first four derivatives of

\[
e^x.
\]
\end{exercise}

\begin{exercise}[3]
Find the first four derivatives of

\[
\sin x.
\]
\end{exercise}

\begin{exercise}[3]
Find

\[
\frac{d^2y}{dx^2}
\]

for

\[
y=x^4.
\]
\end{exercise}

\begin{exercise}[3]
For

\[
x=t^2+1,
\qquad
y=t^3-t,
\]

find

\[
\frac{dy}{dx}.
\]
\end{exercise}

\begin{exercise}[3]
For the same parametric curve, find

\[
\frac{d^2y}{dx^2}.
\]
\end{exercise}

\begin{exercise}[3]
For

\[
r=2\cos\theta,
\]

compute

\[
\frac{dy}{dx}
\]

in polar form.
\end{exercise}

\begin{exercise}[2]
Find the gradient of

\[
f(x,y)=x^2+3xy+y^2.
\]
\end{exercise}

\begin{exercise}[2]
Find the directional derivative of a given function in a specified unit
direction.
\end{exercise}

\begin{exercise}[2]
Find the Jacobian of

\[
F(x,y)
=
(x^2y,e^{xy}).
\]
\end{exercise}

\begin{exercise}[2]
Find the Hessian of

\[
f(x,y)
=
x^3+xy^2.
\]
\end{exercise}

\begin{exercise}[3]
Verify equality of mixed partials for a smooth polynomial.
\end{exercise}

\begin{exercise}[3]
Use the multivariable chain rule for

\[
z=x^2+y^2,
\qquad
x=\cos t,
\qquad
y=\sin t.
\]
\end{exercise}

\begin{exercise}[3]
Differentiate

\[
\mathbf{r}(t)
=
(t^2,e^t,\sin t).
\]
\end{exercise}

\begin{exercise}[3]
Given a position vector, compute velocity, speed, and acceleration.
\end{exercise}

\begin{exercise}[3]
Differentiate

\[
\mathbf{u}(t)\cdot\mathbf{v}(t).
\]
\end{exercise}

\begin{exercise}[3]
Differentiate

\[
\mathbf{u}(t)\times\mathbf{v}(t).
\]
\end{exercise}

\begin{exercise}[3]
Prove

\[
\frac{d}{dt}
\|\mathbf{x}(t)\|^2
=
2\mathbf{x}\cdot\mathbf{x}'.
\]
\end{exercise}

\begin{exercise}[3]
Derive

\[
\frac{d}{dt}
A^{-1}
=
-A^{-1}A'A^{-1}.
\]
\end{exercise}

\begin{exercise}[3]
Differentiate

\[
f(\mathbf{x})
=
\mathbf{a}^T\mathbf{x}.
\]
\end{exercise}

\begin{exercise}[3]
Find the gradient and Hessian of

\[
f(\mathbf{x})
=
\mathbf{x}^TA\mathbf{x}
\]

when \(A\) is symmetric.
\end{exercise}

\begin{exercise}[3]
Find the gradient of

\[
f(\mathbf{x})
=
\frac12
\|A\mathbf{x}-\mathbf{b}\|^2.
\]
\end{exercise}

\begin{exercise}[3]
Differentiate the logistic function

\[
\sigma(x)
=
\frac1{1+e^{-x}}.
\]
\end{exercise}

\begin{exercise}[3]
Show that

\[
\sigma'(x)
=
\sigma(x)(1-\sigma(x)).
\]
\end{exercise}

\begin{exercise}[3]
Differentiate

\[
e^{-x^2}.
\]
\end{exercise}

\begin{exercise}[3]
Differentiate a normal probability density with respect to \(x\).
\end{exercise}

\begin{exercise}[3]
Differentiate

\[
e^{ax}\sin bx.
\]
\end{exercise}

\begin{exercise}[3]
Differentiate

\[
e^{ax}\cos bx.
\]
\end{exercise}

\begin{exercise}[4]
Derive the product rule from the limit definition.
\end{exercise}

\begin{exercise}[4]
Derive the quotient rule from the product rule.
\end{exercise}

\begin{exercise}[4]
Prove the chain rule under suitable differentiability assumptions.
\end{exercise}

\begin{exercise}[4]
Derive the derivative of \(\ln x\) from the inverse function theorem.
\end{exercise}

\begin{exercise}[4]
Derive the derivative of \(\arcsin x\) by implicit differentiation.
\end{exercise}

\begin{exercise}[4]
Derive the derivative of \(\arctan x\).
\end{exercise}

\begin{exercise}[4]
Derive the derivative formula for an inverse function.
\end{exercise}

\begin{exercise}[4]
Derive the derivative of \(x^x\) using logarithmic differentiation.
\end{exercise}

\begin{exercise}[4]
Prove the multivariable chain rule in Jacobian notation.
\end{exercise}

\begin{exercise}[4]
Derive Jacobi's formula for the derivative of a determinant.
\end{exercise}

\begin{exercise}[4]
Derive the gradient of a quadratic form.
\end{exercise}

%%%%%%%%%%%%%%%%%%%%%%%%%%%%%%%%%%%%%%%%%%%%%%%%%%%%%%%%%%%%
\subsection{Conceptual Summary}
\label{subsec:derivative-tables-summary}
%%%%%%%%%%%%%%%%%%%%%%%%%%%%%%%%%%%%%%%%%%%%%%%%%%%%%%%%%%%%

The derivative is defined by

\[
\boxed{
f'(x)
=
\lim_{h\to0}
\frac{
f(x+h)-f(x)
}{
h
}.
}
\]

The central differentiation rules are

\[
\boxed{
(f+g)'=f'+g',
}
\]

\[
\boxed{
(fg)'=f'g+fg',
}
\]

\[
\boxed{
\left(
\frac fg
\right)'
=
\frac{
f'g-fg'
}{
g^2
},
}
\]

and

\[
\boxed{
(f\circ g)'
=
(f'\circ g)g'.
}
\]

Standard elementary derivatives include

\[
\boxed{
(x^n)'=nx^{n-1},
}
\]

\[
\boxed{
(e^x)'=e^x,
}
\]

\[
\boxed{
(\ln x)'=\frac1x,
}
\]

\[
\boxed{
(\sin x)'=\cos x,
}
\]

and

\[
\boxed{
(\cos x)'=-\sin x.
}
\]

Inverse, hyperbolic, parametric, implicit, vector, and multivariable
derivatives extend the same basic idea of local linear approximation.

The central practical principle is

\[
\boxed{
\text{identify the structure first, then apply the corresponding rule}.
}
\]

%%%%%%%%%%%%%%%%%%%%%%%%%%%%%%%%%%%%%%%%%%%%%%%%%%%%%%%%%%%%
\subsection{Section Connections}
\label{subsec:derivative-tables-connections}
%%%%%%%%%%%%%%%%%%%%%%%%%%%%%%%%%%%%%%%%%%%%%%%%%%%%%%%%%%%%

\chapterconnection{
Derivative Tables consolidates differentiation formulas used throughout
Calculus, Analysis, Differential Equations, Optimization, Probability,
Statistics, Geometry, Physics, Continuum Mechanics, Control Theory, and
Mathematical Modeling. The power, product, quotient, and chain rules support
ordinary differentiation; inverse and implicit derivatives connect with
Algebra and the Inverse Function Theorem; gradients, Jacobians, and Hessians
connect with Multivariable Analysis and Optimization; vector derivatives
connect with Geometry and Mechanics; and matrix derivatives support numerical
optimization and applied modeling. The next reference section, Integral
Tables, collects standard antiderivatives, substitution patterns,
trigonometric integrals, inverse-trigonometric forms, hyperbolic forms,
rational-function patterns, Gaussian-type integrals, and common definite
integrals.
}

%%%%%%%%%%%%%%%%%%%%%%%%%%%%%%%%%%%%%%%%%%%%%%%%%%%%%%%%%%%%
% SECTION SOURCE TRACKING
%%%%%%%%%%%%%%%%%%%%%%%%%%%%%%%%%%%%%%%%%%%%%%%%%%%%%%%%%%%%

%------------------------------------------------------------
% 14.9 Derivative Tables
%------------------------------------------------------------
%
% Source basis:
%   - Standard single-variable and multivariable calculus.
%   - Standard derivative identities from analysis, linear algebra,
%     probability, optimization, and applied mathematics.
%
% Used for:
%   - derivative notation;
%   - derivative definition;
%   - rate of change;
%   - constant rule;
%   - power rule;
%   - sum and difference rules;
%   - product and quotient rules;
%   - chain rule;
%   - polynomial derivatives;
%   - radicals;
%   - exponential functions;
%   - logarithmic functions;
%   - trigonometric functions;
%   - inverse trigonometric functions;
%   - hyperbolic functions;
%   - inverse hyperbolic functions;
%   - absolute value;
%   - logarithmic differentiation;
%   - variable powers;
%   - implicit differentiation;
%   - higher derivatives;
%   - inverse-function derivatives;
%   - parametric derivatives;
%   - polar derivatives;
%   - vector-valued derivatives;
%   - velocity and acceleration;
%   - partial derivatives;
%   - gradients;
%   - directional derivatives;
%   - Jacobians;
%   - Hessians;
%   - mixed partial derivatives;
%   - multivariable chain rule;
%   - vector product derivatives;
%   - matrix inverse derivative;
%   - determinant derivative;
%   - log-determinant derivative;
%   - quadratic-form gradients;
%   - least-squares gradients;
%   - logistic derivative;
%   - Gaussian derivatives;
%   - common composite derivatives;
%   - common errors;
%   - exercises.
%
% Notes:
%   - Integral Tables is developed next in Section 14.10.
%   - Series Expansions follows in Section 14.11.
%
%%%%%%%%%%%%%%%%%%%%%%%%%%%%%%%%%%%%%%%%%%%%%%%%%%%%%%%%%%%%